\begin{abstract}

大多模态模型（LMM）已在涉及多种模态的理解与生成任务中展现出卓越能力。尽管这些模型可以接收灵活组合的输入数据，其训练效率仍受两个主要问题影响：异构模型架构造成的流水线阶段不均衡，以及多模态数据多样性所引起的训练数据动态性。

本文提出 \sys{}，一个面向 LMM 训练的动态、模态感知流水线调度框架。\sys{} 通过两项关键技术应对\emph{动态不均衡}挑战：
(1) 将不同模态的计算分离到专用的\emph{流水线段}中，从而在一组连续阶段内部平衡工作负载；
(2) 将输入数据动态拆分为粒度更细、特定于模态的\emph{子微批次}，从而在这些流水线段之间平衡工作负载。
训练期间，\sys{} 利用空闲 CPU 资源异步生成流水线调度，无需停顿训练过程即可针对每个输入批次动态定制阶段执行。
我们在五种多样化的 LMM 上验证了 \sys{}；这些模型的参数规模从 12B 到 94B，涵盖视觉语言模型和扩散模型。
实验结果表明，与最先进的系统相比，本系统可将吞吐量最高提升 97.3\%，展现出对波动的多模态训练工作负载的强适应能力。

\end{abstract}
